\documentclass[answers,12pt,addpoints]{exam} 
\usepackage{import}

\import{C:/Users/prana/OneDrive/Desktop/MathNotes}{style.tex}

% Header
\newcommand{\name}{Pranav Tikkawar}
\newcommand{\course}{Spatial Stats}
\newcommand{\assignment}{Homework 6.1.1}
\author{\name}
\title{\course \ - \assignment}

\begin{document}
\maketitle

\newpage
\begin{questions}
\question I see an new benefit of reparametrization to ensure that our bayesian priors are orthogonal. In an arbitrary function with 3 or more parameters does these exist a mutually orthogonal set. I feel like this would be true, and it would require $n$ invertable functions for $n$ variates ${f_i}_{i=1^n}$ such that $\phi_i = f_i(\Theta)$ and $\cov(\phi_i, \phi_j) = 0$ and we can generate them by some PCA method or a "stochastic" gram-schmidt orthogonalization. 

\question the notion of a process "rougher" than brownian motion is brought up, and I wanted to just confirm that this is in the path sense (how differentiable), not quadratic variation sense. I remember earlier in the year (on the snow day) the notion was brought up. 

Honestly the reason why I am getting confused is due to the fact that i know that for a R.V with finite second moment, the Q.V should be 0, but if we have a GP which is rougher (and by my assumption larger QV) than Brownian motion, it should have infinite second moment, but being a GP that is a contradiction.

\question The RMSE for $\hat \phi_2$ does not change much at all between models which was interesting. I feel like the reason might have to due with the magnitude of the $\eta$ and $\beta$ parameters as with smaller parameters for the eccentricity the less we need to adjust the estimates from a fixed model.




\end{questions}

\end{document}